\documentclass[11pt,letterpaper]{article}
\usepackage[margin=0.85in]{geometry}
\usepackage{amsmath,amssymb,booktabs,graphicx}
\usepackage[T1]{fontenc}
\usepackage{lmodern}
\usepackage[colorlinks=true,urlcolor=blue]{hyperref}
\setlength{\parindent}{0pt}
\setlength{\parskip}{7pt}
\begin{document}
\begin{center}
{\Large\bfseries Training 2-bit weights with almost no optimizer state}\\[5pt]
{\small Draft, 2026-11-10. Rotation rounding vs.\ flip, GPEF and quantized-moment baselines}
\end{center}

\section*{1. Question}

A model whose weights are stored on a 2-bit grid is usually trained with a 32-bit shadow copy and 64 bits of
Adam moments per weight, about 100 bits of training state for a 2-bit deployed weight. We ask how much accuracy each
method buys per bit of training state, especially below 4 bits per weight, where a method can afford at most two
extra bits beyond the deployed code. We compare one momentum-free method family (SRQ: a normalized gradient step
with stochastic or rotation rounding onto a sub-grid) with a teammate's flip optimizer (quantized momentum), a
teammate's GPEF (error feedback with learned group codebooks), and Adam/Muon with stochastic rounding and quantized
moments.

\section*{2. Setting}

\textbf{Format.} Every quantized matrix $W$ deploys codes $K_{ij}\in\{-1.5,-0.5,0.5,1.5\}$ and one scale
$\alpha$ per matrix, $W=\alpha K$. GPEF and AdamW-STE instead deploy a learned codebook $\{-a,-b,b,a\}$ for every
group of 32 weights (2 extra bits per weight in fp32 codebook storage). Embeddings, head and norms are fp32 AdamW in
all methods.

\textbf{State accounting.} \emph{Training state per weight} counts everything kept per quantized weight across
steps: the 2-bit code plus optimizer state. The transient gradient is not counted.

\textbf{Models and data.} Byte-level transformer language models from the flip harness, trained from scratch:
2.65M quantized weights on WikiText-2 for 2{,}500 steps (pilot), and 10.6M on WikiText-103 for 8{,}000 steps
(batch $64\times256$).

\textbf{Protocol.} For each method: a seed-0 learning-rate grid of three values spaced $3\times$, extended by
$3\times$ at an edge optimum, then seeds 1 and 2 at the seed-0 best. Reported numbers are validation bits per byte
(mean of 3 seeds). Test data is never read. All final learning rates are interior to the searched grid.

\section*{3. Methods}

\textbf{SRQ (srq\_sr@b$B$, srq\_rot@b$B$; $2+B$ bits/weight).} Each weight has a position $x=K+f$, where $f$ is a
remainder on the sub-grid $2^{-B}\mathbb Z\cap[-\tfrac12,\tfrac12)$; only $K$ is deployed. One step, with no
momentum, is
\[
\tilde x=\operatorname{clip}\!\left(x-\eta\,\frac{g}{\operatorname{rms}(g)},\,-1.5,\,1.5\right),
\qquad g=\partial L/\partial W,
\]
followed by rounding $\tilde x$ onto the $2^{-B}$ sub-grid; $K$ is the nearest code and $f$ the remainder. With
$B=0$ the position is the code itself (2 bits/weight). \emph{SR} rounds up with probability equal to the fractional
part. \emph{Rotation} rounds up when the fractional part exceeds the threshold $\operatorname{frac}(u_0+t\varphi)$,
with $\varphi=(\sqrt5-1)/2$ and $u_0$ regenerated per entry from a per-matrix seed, so it adds no state. Over
consecutive steps the thresholds are equidistributed with small gaps, which lowers the variance of accumulated
rounding error relative to independent draws.

\textbf{Flip (flip@m$B$; $2+B+\approx0.1$ bits/weight).} Teammate optimizer: $B$-bit quantized momentum per
weight, and each step moves the top-$k$ weights by one code level under row and column caps. 1.58 = ternary
momentum; ``2g'' = 2-bit momentum with group scales.

\textbf{GPEF (gpef@e$E$; $2+E+2$ bits/weight).} Teammate method: an rms-normalized step in units of the group
codebook, two least-squares codebook/symbol passes per step (each symbol set to the nearest level), and $E$-bit
deterministic error feedback ($E=0$: none). To go below 4 bits/weight, gpef@e0g$G$ shares one codebook among $G$
contiguous weights instead of 32 ($2+64/G$ bits/weight: $G=64$, 128, 256 give 3, 2.5, 2.25); $G=32$ reproduces
gpef@e0 exactly. As a format-matched check of the rounding rule, gpef\_rot@e0 and gpef\_sr@e0 keep GPEF's step and
codebook fit but choose each symbol by rotation or stochastic rounding between the two bracketing levels.

\textbf{References.} Adam + SR with both moments quantized to $B$ bits (adam\_sr@m$B$; $2+2B$); Muon + SR with
$B$-bit momentum (muon\_sr@m$B$; $2+B$); Muon + QAT with fp32 shadow (66); AdamW-STE (teammate, 102).

\section*{4. Results}

\begin{table}[h]\centering\small
\begin{tabular}{lrrr}
\toprule
method & bits/weight & 2.65M & 10.6M\\
\midrule
srq\_rot@b0 & 2 & \textbf{2.519} & \textbf{2.334}\\
srq\_sr@b0 & 2 & 2.570 & 2.391\\
gpef@e0g256 & 2.25 & \textbf{2.450} & \textbf{2.221}\\
gpef@e0g128 & 2.5 & 2.377 & 2.133\\
srq\_rot@b1 & 3 & 2.332 & --\\
srq\_sr@b1 & 3 & 2.383 & --\\
gpef@e0g64 & 3 & \textbf{2.260} & \textbf{2.019}\\
flip@m1.58 & 3.7 & 2.416 & 2.204\\
srq\_rot@b2 & 4 & 2.168 & 1.941\\
srq\_sr@b2 & 4 & 2.228 & 2.055\\
gpef@e0 (group-32 codebooks) & 4 & \textbf{2.135} & \textbf{1.889}\\
gpef\_rot@e0 (rotation symbols) & 4 & 2.505 & 2.261\\
gpef\_sr@e0 (SR symbols) & 4 & 2.555 & 2.316\\
flip@m2 & 4.1 & 2.327 & 2.124\\
muon\_sr@m2 & 4.1 & 2.641 & 2.545\\
flip@m2g & 4.25 & 2.242 & 2.058\\
\midrule
srq\_rot@b4 & 6 & 1.959 & 1.689\\
srq\_sr@b4 & 6 & 1.976 & 1.721\\
flip@m4 & 6.1 & \textbf{1.873} & \textbf{1.676}\\
adam\_sr@m2 & 6.25 & 2.606 & 2.524\\
srq\_rot@b6 & 8 & 1.934 & \textbf{1.635}\\
srq\_sr@b6 & 8 & 1.938 & 1.636\\
gpef@e4 & 8 & 2.054 & 1.758\\
adam\_sr@m4 & 10.25 & 2.111 & 1.894\\
gpef@e8 & 12 & 1.879 & 1.607\\
\midrule
muon\_qat (fp32 shadow) & 66 & 1.745 & 1.530\\
adamw\_ste & 102 & 1.826 & --\\
\bottomrule
\end{tabular}
\caption{Validation bits per byte, mean of 3 seeds (seed sd mostly .002--.02; srq\_rot@b2 at 10.6M .034).
Pilot-only points (flip@m8 1.853 at 10.1, srq@b8 1.906/1.918 at 10, adam\_sr@m8 2.074, Adam+SR with fp32
moments 2.067) were not rerun at 10.6M.}
\end{table}

\begin{figure}[h]\centering
\includegraphics[width=\linewidth]{bits_vs_bpb_10M.png}
\caption{10.6M model. Lines: mean of 3 seeds at the best learning rate; dots: individual seeds.}
\end{figure}

\section*{5. Findings}

\textbf{Below 4.5 bits per weight, the teammate's GPEF is the best method at every budget from 2.25 to 4
bits/weight.} With one learned codebook per 256 weights (2.25 bits/weight) it reaches 2.221 at 10.6M, .11 better
than rotation rounding with no state at all (2.334 at 2 bits/weight). At 3 bits/weight gpef@e0g64 beats
srq\_rot@b1 by .07 (2.65M) and flip@m1.58 at 3.7 bits by .19 (10.6M). At 4 bits/weight gpef@e0 beats srq\_rot@b2
by .03 and .05. Its curve is smooth in group size: each doubling of the group costs about .1 bpb.

\textbf{On the fixed per-matrix format, rotation rounding is the best method below 4.5 bits per weight.} At 4
bits/weight it beats flip's best low-bit variant (2-bit grouped momentum, 4.25 bits) by .07 at 2.65M and .12 at
10.6M, and plain SR by .06 and .11. At 2 bits/weight, with no state beyond the code, it trains to 2.33. Flip does
not fail below 4-bit momentum, but it degrades by .37--.45 bpb when its momentum drops from 4 to 2 bits.

\textbf{Above 6 bits per weight the methods converge.} Flip with 4-bit momentum wins at 6 bits/weight, by .09 at
2.65M but only .013 at 10.6M. At 8--10 bits/weight rotation and SR are within seed noise (1.635 vs 1.636 at
10.6M), as expected once the sub-grid is fine enough that rounding variance is negligible.

\textbf{GPEF's advantage comes from its learned codebooks with nearest-level symbols, not from its rounding rule.}
GPEF@e0 has no error feedback and no momentum. Its learned codebooks also change the deployed model: per-group
levels instead of our 2 bits plus one scale per matrix. In the format-matched check, choosing GPEF's symbols by
rotation or stochastic rounding instead of nearest-level makes it much worse (10.6M: 2.261 and 2.316 against 1.889).
Rotation still beats SR there by .05 at both sizes, so the rounding-rule ranking holds on a non-uniform format. Our
reading, not yet tested, is that the fp32 codebook levels already integrate the small updates, so randomizing the
symbols only adds noise. With error feedback GPEF spends extra bits less efficiently than SRQ (gpef@e4 1.758 vs
srq\_rot@b6 1.635 at 8 bits/weight).

\textbf{Adam- and Muon-direction updates with stochastic rounding are poor at any budget.} Adam + SR stays near 2.07
at 2.65M even with fp32 moments, and collapses with 2-bit moments; Muon + SR at 4--6 bits/weight is .4--.6 worse
than flip and SRQ. The limit is the combination of these directions with rounding onto 2-bit codes, not moment
precision.

\textbf{Gap to full-precision training.} At 10.6M the best method within 8 bits/weight (srq\_rot@b6, 1.635) is .10
from Muon + QAT with an fp32 shadow (1.530, 66 bits/weight); within 6 bits/weight flip@m4 is .15 away.

\section*{6. Limitations and next steps}

\begin{itemize}\setlength\itemsep{2pt}
\item Two sizes (2.65M, 10.6M) and one data family. 25.2M (about 3 h per run) not yet run.
\item GPEF and SRQ differ in deployed format as well as in update rule. The format-matched check shows that
  stochastic symbols hurt on GPEF's format; the converse test, GPEF's nearest-level update on the per-matrix grid,
  has not been run, so how much of GPEF's lead is format alone is open.
\item srq@b1 was run at 2.65M only.
\item The per-row scale refit (srqrr) hurt at 2.65M and was dropped; at 10.6M it has seed 0 only.
\item Flip@m8, srq@b8, AdamW-STE and Adam + SR with fp32 moments were run at 2.65M only.
\end{itemize}

\end{document}
